\chapter*{Streszczenie}
\addcontentsline{toc}{chapter}{Streszczenie}
\thispagestyle{plain}

Celem tego projektu by\l{}o stworzenie aplikacji desktopowej, kt\'{o}ra ma u\l{}atwi\'{c}
prac\k{e} o\'{s}rodka szkolenia kierowc\'{o}w. Aplikacja powsta\l{}a w~j\k{e}zyku C\#
z~wykorzystaniem technologii WPF i~wzorca MVVM, a~dane przechowywane s\k{a} w~lokalnej
bazie SQL~Server obs\l{}ugiwanej przez Entity~Framework~Core.

W~aplikacji zosta\l{}y zaimplementowane trzy role u\.{z}ytkownik\'{o}w. Administrator
mo\.{z}e dodawa\'{c}, edytowa\'{c} i~usuwa\'{c} kursant\'{o}w, instruktor\'{o}w oraz
pojazdy, a~tak\.{z}e zarz\k{a}dza\'{c} wszystkimi rezerwacjami. Instruktor widzi sw\'{o}j
grafik na dany tydzie\'{n} i~mo\.{z}e anulowa\'{c} w\l{}asne jazdy. Kursant mo\.{z}e sam
rezerwowa\'{c} dost\k{e}pne terminy jazd przez interaktywny kalendarz.

Projekt zawiera mechanizmy walidacji danych, bezpieczne przechowywanie hase\l{}
oraz ochron\k{e} integralno\'{s}ci danych w~bazie.
